\section{Part 3：组合并行与真实模型配置}

单个 primitive 解决单一瓶颈，真实训练则必须把多个维度映射到硬件层次。本部分从“先 fit、再优化吞吐”的顺序出发，解释为什么高频 TP/EP 通常留在单节点快链路，PP 跨节点切层，DP/FSDP 再在更大范围复制或分片训练状态。

\subsection{3D/4D parallelism 的经验规则}

本节把组合问题写成一组可执行规则，而不是背诵某个模型的数字。首先根据参数与 activation 账本确定最小 TP/PP/EP/CP 规模让模型可训练；然后用剩余设备扩大 DP、增加 microbatch、重叠通信，并通过 profiling 验证每一维是否真的缓解了对应瓶颈。

\begin{figure}[H]
\centering
\includegraphics[width=0.86\textwidth]{slide-057.jpg}
\caption{Slide 57：3D/4D parallelism 的简单经验规则：先让模型 fit，再优化吞吐。}
\figsource{CS336 Spring 2026 Lecture 8 官方幻灯片。}
\end{figure}

\begin{importantbox}{读图：Slide 57 的规则顺序}
第一，模型放不下时，优先用 TP/EP 到单机 GPU 上限，再用 PP 跨机器，或者使用 ZeRO-3/FSDP。第二，模型能放下后，用 DP/FSDP 增加吞吐。第三，TP/EP 通常限制在高速节点内，PP 可跨节点，DP 可扩到更大范围。这个顺序体现了“先解决容量，再解决吞吐”的工程逻辑。
\end{importantbox}


\singleslide{slides-images/slide-058.jpg}{Megatron 当前推荐：基础 DP/TP/PP/SP 之上，再按需要加入 CP、EP 与更细重叠优化。}{58}

Slide 58 将经验规则映射到具体框架选项。“Basics”先解决模型是否 fit 与主要吞吐，“extras”才处理长上下文、MoE、通信重叠与高级调度。配置顺序很重要：未测清基础瓶颈就叠加所有维度，会让 group construction 和性能归因变得困难。

\begin{knowledgebox}{讲义补充：源 Slide 58 的 Megatron 建议怎么读}
Megatron 的建议不是固定答案，而是一组经过实践验证的起点。基础配置通常先选 TP、PP、DP；更复杂模型再引入 EP、CP、sequence parallel、activation recomputation 和 communication overlap。实际配置要根据模型结构、硬件域和 profiling 调整。
\end{knowledgebox}

\begin{figure}[H]
\centering
\includegraphics[width=0.86\textwidth]{slide-059.jpg}
\caption{Slide 59：Narayanan 2021 的 scaling strategies，TP 先到 8，PP 逐步增加，DP 逐步减少。}
\figsource{CS336 Spring 2026 Lecture 8 官方幻灯片。}
\end{figure}

\begin{knowledgebox}{读表：Slide 59 的趋势}
表中 TP 很快达到 8 并封顶，通常对应单节点 8 GPU 的高速互连域。模型继续变大时，PP 增加以容纳更多层和参数；总 GPU 数固定或增长时，DP size 反而下降，因为更多 GPU 被用来放模型而不是复制模型处理更多 batch。
\end{knowledgebox}

\begin{figure}[H]
\centering
\includegraphics[width=0.86\textwidth]{slide-060.jpg}
\caption{Slide 60：谨慎的 3D parallelism 可以获得接近线性的 gains 和平坦 utilization。}
\figsource{CS336 Spring 2026 Lecture 8 官方幻灯片。}
\end{figure}

\begin{knowledgebox}{读图：Slide 60 的“linear gains”有前提}
接近线性扩展来自把高频通信放在快链路、把模型切分均衡、把 bubble 控制住、让 DP 规模不超过有效 batch。它不是并行维度越多越好，而是每个维度都服务于明确瓶颈。
\end{knowledgebox}


\singleslide{slides-images/slide-061.jpg}{跨 64 台机器的经验甜点：单节点内 TP=8，节点间用其他并行维度扩展。}{61}

Slide 61 的 8×8 配置把通信频率与硬件域对齐。TP 每层都通信，应留在单节点 8-GPU 高速域；跨节点更适合较低频的 pipeline 或 data parallel。TP=8 不是普适常数，而是当前常见节点拓扑下的局部最优。

\begin{knowledgebox}{讲义补充：源 Slide 61 为什么 TP=8 常出现}
8 常对应单节点 GPU 数或一个高速互连域大小。TP 再扩大可能跨节点，引入高频跨节点通信；TP 太小又可能单层放不下或每卡计算太重。因此 TP=8 是许多 GPU 集群上的硬件甜点，而不是数学常数。
\end{knowledgebox}

\begin{figure}[H]
\centering
\includegraphics[width=0.86\textwidth]{slide-062.jpg}
\caption{Slide 62：activation recomputation 通过省显存允许更大 batch，可能反而提高吞吐。}
\figsource{CS336 Spring 2026 Lecture 8 官方幻灯片。}
\end{figure}

\begin{knowledgebox}{读图：Slide 62 为什么 recomputation 可能“pay for itself”}
Activation recomputation 或 checkpointing 用额外 forward 计算换显存。如果省下的显存允许更大 microbatch/global batch，GPU 利用率可能提高，pipeline bubble 可能下降，整体吞吐反而变好。因此它不是简单“慢一点换省显存”，而是可能改变可行 batch regime。
\end{knowledgebox}

\subsection{近期模型配置案例}

最后用公开模型配置检验这些规则。不同模型的 DP/TP/PP/EP/CP 数字并不是互相矛盾，而是架构、参数规模、上下文长度、节点拓扑和训练阶段不同的结果；本节重点不是记住配置表，而是从每个案例反推“哪个状态放不下、哪个通信必须留在快链路”。


\singleslide{slides-images/slide-063.jpg}{近期模型案例入口：7B 级 Dolma 主要用 FSDP，说明小模型不必启用全部并行维度。}{63}

Dolma 案例提供了复杂度下界：当模型能在少量节点内借助 FSDP 放下时，TP/PP/EP 的收益可能不抵实现成本。并行方案应从最简单满足 memory/throughput 的配置开始，再用 profiling 证明需要新增维度。

\begin{knowledgebox}{讲义补充：源 Slide 63 的意义}
小到 7B 量级的模型在现代多卡节点上可能主要靠 FSDP/ZeRO-3 解决显存和吞吐，不一定需要复杂 TP/PP/EP。并行复杂度应与模型规模匹配，过早引入所有并行维度会增加工程风险。
\end{knowledgebox}

\begin{figure}[H]
\centering
\includegraphics[width=0.86\textwidth]{slide-064.jpg}
\caption{Slide 64：DeepSeek 使用 ZeRO stage 1，并组合 TP、SP、PP、EP 等策略。}
\figsource{CS336 Spring 2026 Lecture 8 官方幻灯片。}
\end{figure}

\begin{knowledgebox}{读图：Slide 64 的 DeepSeek 配置}
DeepSeek V3 中 PP、EP、ZeRO-1、all-to-all overlap 等同时出现。MoE 模型的核心压力来自 experts 和 token routing，因此 EP 规模很大；attention/activation 相关部分仍需要 TP/SP/PP。这里体现了异构模型结构需要异构并行策略。
\end{knowledgebox}


\singleslide{slides-images/slide-065.jpg}{Yi 系列配置演进：dense Yi 使用 ZeRO-1 + TP + PP，Yi-Lightning 转向 expert parallel。}{65}

Slide 65 是架构变化驱动并行变化的直接例子。Dense MLP 适合 tensor/pipeline slicing；模型改成 MoE 后，参数和计算沿 experts 自然分组，EP 能让每个 token 只访问激活专家。系统配置必须追随计算图，而不是固定继承上一代模型。

\begin{knowledgebox}{讲义补充：源 Slide 65 说明并行策略会随架构变化}
Dense 模型常依赖 TP/PP/DP；当模型走向 MoE，expert parallelism 会替代一部分 dense tensor parallel。策略变化不是追潮流，而是模型计算形态从“一个大 dense MLP”变成“多个 experts 加 routing”。
\end{knowledgebox}

\begin{figure}[H]
\centering
\includegraphics[width=0.86\textwidth]{slide-066.jpg}
\caption{Slide 66：Llama 3 405B 的多阶段训练配置，包括 Stage 1/2/3 和 long-context。}
\figsource{CS336 Spring 2026 Lecture 8 官方幻灯片。}
\end{figure}

\begin{knowledgebox}{读图：Slide 66 的阶段化训练}
同一个模型在不同阶段可能使用不同并行和 batch 配置。早期小 batch、主预训练、长上下文扩展的瓶颈不同：有时是参数显存，有时是吞吐，有时是 KV/activation 和 sequence length。因此并行配置是训练计划的一部分，不是一次性固定。
\end{knowledgebox}

\begin{figure}[H]
\centering
\includegraphics[width=0.86\textwidth]{slide-067.jpg}
\caption{Slide 67：Llama 3 405B 规模下 GPU failures 很常见，可靠性成为训练系统问题。}
\figsource{CS336 Spring 2026 Lecture 8 官方幻灯片。}
\end{figure}

\begin{warningbox}{读图：Slide 67 的系统含义}
GPU 数越多，单个硬件故障的期望频率越高。并行训练不只要跑得快，还要 checkpoint、恢复、容错、监控和动态剔除故障节点。大型训练的系统工程包括可靠性，而不仅是并行算法。
\end{warningbox}


\singleslide{slides-images/slide-068.jpg}{Gemma 2 的 dense 配置：ZeRO-3、model parallel（TP+SP）与 data parallel。}{68}

Gemma 2 的 2B/9B/27B 版本采用较传统 dense 路线：ZeRO-3 分训练状态，TP+SP 同时切矩阵与 activation，剩余 GPU 用 DP 提升吞吐。没有 PP/EP/CP，反映模型深度、上下文和单节点容量尚不需要所有维度。

\begin{knowledgebox}{讲义补充：源 Slide 68 的 Gemma 2 配置}
Gemma 2 的例子说明，即使没有 MoE，也常组合 FSDP/ZeRO-3、TP、SP、DP。参数、activation 和吞吐各自需要不同维度来解决，尤其当模型有多个尺寸版本时，并行策略也会随模型大小调整。
\end{knowledgebox}


\singleslide{slides-images/slide-069.jpg}{Mixtral 8×22B：TP/PP/CP/EP = 4/4/1/8，并以 DP 补齐总 GPU 数。}{69}

Mixtral 同时需要 TP 处理 attention/dense 部分、EP 切 experts、PP 切层，CP=1 表示上下文不是主要瓶颈。配置乘积必须与 world size 一致，剩余 factor 才能作为 DP；这也是阅读公开并行数字时最基本的自洽检查。

\begin{knowledgebox}{讲义补充：源 Slide 69 的 Mixtral MoE 配置}
Mixtral 是 MoE 模型，因此 EP 很自然出现；TP/PP/CP 处理 dense attention、深度切分和长上下文；DP 补充吞吐。MoE 模型的并行配置通常更高维，因为 MLP experts、attention、context 和 batch 的瓶颈不同。
\end{knowledgebox}


\singleslide{slides-images/slide-070.jpg}{Nemotron 3 Super 长上下文扩展：TP/CP/EP = 2/64/64，sequence 轴成为主要分片对象。}{70}

CP=64 表示单条长序列的 attention/KV/activation 已无法只靠参数并行解决，需要大量 ranks 共同处理 context。与此同时 EP=64 处理 MoE experts，两个大并行 group 可能需要不同通信模式与拓扑映射；长上下文 MoE 是典型的多维耦合系统。

\begin{knowledgebox}{讲义补充：源 Slide 70 的 CP=64 为什么醒目}
长上下文扩展会让 attention 和 KV/activation 压力显著上升，因此 context parallel 可能成为主维度之一。CP 数字很大说明序列维度本身已经成为核心扩展对象，而不只是把参数切开。
\end{knowledgebox}


\singleslide{slides-images/slide-071.jpg}{Qwen3 MoE：小规模优先节点内 EP，更大模型使用 2/8/32 等分层并行配置。}{71}

Qwen3 的 A22B/A3B 表示每 token 激活参数远小于总参数，因此 expert storage、routing 与 load balance 比 dense FLOPs 更关键。节点内最多 8-way EP 可以利用快互连；规模继续扩大时，再叠加其他维度。并行数字反映的是 activated computation 与网络域，而非总参数量本身。

\begin{knowledgebox}{讲义补充：源 Slide 71 的 Qwen 配置}
Qwen 3 的 MoE 配置说明 expert parallel 可以成为主要扩展方式。A22B/A3B 代表 activated parameters 少于总参数，系统瓶颈不只是“总参数存储”，还包括每 token 激活专家数量、路由通信和 expert load balance。
\end{knowledgebox}

\begin{figure}[H]
\centering
\includegraphics[width=0.9\textwidth]{slide-072.jpg}
\caption{Slide 72：模型并行 overview table，横向比较 DeepSeek、Yi、Llama3、Gemma2、Mixtral、Nemotron、Qwen 等配置。}
\figsource{CS336 Spring 2026 Lecture 8 官方幻灯片。}
\end{figure}

\begin{importantbox}{读表：Slide 72 不要把公开配置当唯一真相}
这张表最有价值的是比较维度：DP、TP/SP、EP、PP、CP 如何随模型类型变化。Dense 模型更依赖 TP/SP/PP/DP；MoE 模型显著引入 EP；长上下文模型提高 CP。公开配置可能不完整或随训练阶段变化，因此它们应作为推理线索，而不是可直接复制的 recipe。
\end{importantbox}

\begin{figure}[H]
\centering
\includegraphics[width=0.86\textwidth]{slide-073.jpg}
\caption{Slide 73：全讲 recap，没有单一并行方案，通常需要组合多种 approach，并遵循可解释规则。}
\figsource{CS336 Spring 2026 Lecture 8 官方幻灯片。}
\end{figure}

\begin{importantbox}{读图：Slide 73 的最终结论}
超过某个规模后，multi-GPU、multi-node parallelism 是必要条件；没有单一方案能解决所有问题；可解释的经验规则通常先问模型是否 fit，再问吞吐如何扩展，最后按硬件域安排 TP/EP/PP/DP/CP。真正可靠的配置来自账本、benchmark 和 profiling，而不是名词堆叠。
\end{importantbox}

\subsection{Part 3 小结}

大模型训练的组合并行遵循一个朴素原则：\textbf{把高频通信放在快链路，把大状态按最自然的维度分片，把调度空泡和负载不均控制在可接受范围内}。近期模型配置之所以复杂，是因为现代 LLM 同时有 dense attention、MoE experts、长上下文、深层网络和庞大训练 batch。

\subsection{本章小结}

Lecture 8 把网络拓扑、collectives、状态分片、模型切分和真实模型配置连成一张系统账本。前两部分回答“为什么要并行、有哪些切法”，最后一部分回答“怎样把这些切法按硬件层次组合起来”。

